
\section{1. Introduction}

\subsection{1.1 Motivation}

Transformer inference at long sequence length is dominated by the 
$\mathcal{O}(L^2)$ cost of softmax attention. Two broad strategies have 
emerged to reduce this cost:

1. \textbf{Dense softmax attention} $\mathcal{A}\_\text{dense}(Q, K, V) = 
   \text{softmax}(QK^\top / \sqrt{d}) V$. This preserves the expressive power 
   of exact attention but requires materializing an $L \times L$ matrix.

2. \textbf{Linear attention} $\mathcal{A}\_\text{lin}(Q, K, V) = 
   \phi(Q) (\phi(K)^\top V) / (\phi(Q) \cdot \sum\_j \phi(K\_j))$ for a kernel 
   $\phi$. This reduces complexity to $\mathcal{O}(L)$ but trades quality.

Neither extreme is universally better. On tasks requiring fine-grained 
reasoning (e.g., in-context learning), dense attention wins. On retrieval-like 
tasks with local structure, linear attention matches dense at a fraction of 
the cost.

\textbf{The opportunity.} Rather than commit to one strategy, \textbf{hybrid attention} 
runs both and combines them. If the combination weight is learned per-head 
and per-token, the model can allocate dense capacity where it matters and 
linear capacity elsewhere.

\subsection{1.2 Problem Statement}

Given a pretrained transformer with dense attention, we seek a modified 
module $\mathcal{A}'$ such that:

1. \textbf{Efficiency:} $\mathcal{O}(L)$ memory and time for the linear branch, 
   at most $\mathcal{O}(L^2)$ for the dense branch when needed.
2. \textbf{Quality:} $\Delta\text{Acc} \leq 5\%$ on SST-2.
3. \textbf{Pretrained compatibility:} No retraining of base weights; only LoRA 
   and gating parameters are trainable.
4. \textbf{Architecture portability:} Flat, nested, and causal transformers.

\subsection{1.3 Our Approach: HALO-ReGLA}

We propose \textbf{HALO-ReGLA} (Hardware-Aware Linear Optics \& Rectified Gated 
Linear Attention), which:

1. \textbf{Runs two attention streams in parallel.}
\begin{itemize}
\item Dense: standard softmax attention.
\item Linear: rectified kernel attention using $\phi(x) = \mathrm{elu}(x) + 1$.
\end{itemize}

2. \textbf{Combines them via per-head learnable gates.} The gate for head $h$ at 
   position $i$ is $g_{h,i} = \sigma(w_h^\top x_i + b_h)$, and the output is 
   $g \odot O_\text{dense} + (1 - g) \odot O_\text{lin}$.

3. \textbf{Initializes the gate to favor dense behavior.} We set $b_h = 1.8$ so 
   $\sigma(1.8) \approx 0.86$ at initialization. This ensures the model 
   starts near the pretrained dense behavior.

4. \textbf{Uses a rectified kernel for numerical stability.} The non-negativity of 
   $\phi(x) = \mathrm{elu}(x) + 1$ guarantees the denominator 
   $\phi(Q) \cdot \sum_j \phi(K_j)$ is positive, avoiding division by zero.

5. \textbf{Preserves causality via prefix sums.} For causal models, the linear 
   branch uses a cumulative-sum formulation:
\[
   O_i^\text{lin} = \frac{\phi(Q_i) \sum_{j \le i} \phi(K_j) V_j^\top}{\phi(Q_i) \sum_{j \le i} \phi(K_j)}
\]

\subsection{1.4 Contributions}

\begin{itemize}
\item \textbf{A hybrid attention module.} HALO-ReGLA runs dense and linear streams
\end{itemize}
  in parallel and combines them with per-head learnable gates.
\begin{itemize}
\item \textbf{A numerically stable formulation.} The rectified kernel
\end{itemize}
  $\phi(x) = \mathrm{elu}(x) + 1$ ensures positive denominators in both 
  causal and bidirectional settings.
\begin{itemize}
\item \textbf{A production-hardened implementation.} Supports four architectures
\end{itemize}
  with RoPE, GQA, LoRA recovery, and safe injection.
\begin{itemize}
\item \textbf{A systematic evaluation.} Benchmarked on DistilBERT, RoBERTa, and GPT-2
\end{itemize}
  on SST-2 with ablations over gate bias, kernel choice, and fusion strategy.
\begin{itemize}
\item \textbf{Gate analysis.} We show that learned gates are non-uniform and
\end{itemize}
  systematically depend on layer depth and head index.

\subsection{1.5 Paper Organization}

Section 2 reviews related work. Section 3 provides background on dense and 
linear attention. Section 4 presents HALO-ReGLA. Section 5 describes 
integration with four architectures. Section 6 details the experimental 
setup. Section 7 reports results. Section 8 analyzes the learned gates. 
Section 9 discusses practical guidance. Section 10 lists limitations. 
Section 11 concludes.



\subsection{Notation}
